\section{Le découplage CRT et la non-localité modulaire}
\label{sec:decouplage_crt}

Les Sections~\ref{sec:champ_connexion} et~\ref{sec:holonomie_primitive}
ont établi le statut de l'angle d'holonomie $\theta_p$ comme
invariant primitif et celui des objets classiques $A_\mu,
F_{\mu\nu}, g_{\mu\nu}$ comme dérivés du potentiel scalaire $S$.
Pour relier ce statut à la question expérimentale de la
non-localité, il reste à examiner les corrélations entre canaux
modulaires distincts --- les seules corrélations
``inter-objets'' que la PT autorise à porter une signification
informationnelle non-triviale. Cette section énonce, sans
démonstration détaillée, les deux théorèmes structurels prouvés
dans l'article companion technique~\cite{PT_CRT_DECOUPLING}
ainsi que leur corollaire empirique.

\subsection{Théorème 4.1 : Théorème de découplage géométrique (forme empirique)}
\label{ssec:thm_4_1_fr}

\begin{theoreme}[Théorème de découplage géométrique, forme empirique]
\label{thm:invariance_geom}
\footnote{Vérifié formellement en Lean~4 + Mathlib comme
\texttt{PT.CrtDecoupling.Empirical.empirical\_invariance}
(forme hypothétique),
\texttt{SpectralReduction.empirical\_invariance\_closed}
(forme fermée fini),
\texttt{Phase3.empirical\_invariance\_infinite}
(forme infinie via Birkhoff--Hopf), et
\texttt{Phase4.sieve\_empirical\_invariance}
(clôture spécifique au crible) ;
cf.~\cite[\S~8.5]{PT_CRT_DECOUPLING}.}
Soient $p, q \in \{3, 5, 7\}$ premiers actifs distincts, et soient
$\Phi_p \in \mathcal{A}_p$, $\Phi_q \in \mathcal{A}_q$ des
observables bornées sur les $\sigma$-algèbres modulaires
correspondantes. La fonction
\begin{equation}
  \mathcal{I}_{p,q}(\mu)
  \;:=\;
  I_{\pi_m^{\mathrm{emp}}}\!\bigl(\Phi_p\,;\,\Phi_q \,\big|\, G(\mu)\bigr)
\end{equation}
est constante sur $\mu \in [\mu_c, \infty)$.
\end{theoreme}

La démonstration complète est donnée dans le théorème~6.1
du companion technique \cite{PT_CRT_DECOUPLING}. Elle procède en
quatre étapes : (i) la mesure empirique $\pi_m^{\rm emp}$ de la
trajectoire des résidus est une fonctionnelle du seul champ
$\{g_n\}$, donc indépendante de $\mu$ ; (ii) la géométrie
$G(\mu)$ est mesurable par rapport à la tribu engendrée par le
spectre de la matrice de transfert $T_m$, qui est elle-même
contenue dans la tribu invariante du shift ; (iii) la dynamique du
crible est ergodique par le théorème de Birkhoff--Hopf appliqué
au gap spectral T4, donc la tribu invariante est triviale
modulo $\pi_m^{\rm emp}$, donc $G(\mu)$ est presque sûrement
constante ; (iv) conditionner par une variable presque sûrement
constante ne change pas l'information mutuelle. La conclusion
suit immédiatement.

Le point conceptuel est que la métrique émergente $G(\mu)$ --- y
compris à travers la transition de Hartle--Hawking en $\mu_c
\approx 6{,}97$ et jusqu'au point fixe canonique
$\mustar = 15$ --- ne modifie jamais le contenu informationnel
joint des canaux modulaires actifs. Les corrélations
inter-canaux sont \emph{intrinsèquement} pré-géométriques : elles
existent avant la construction de la métrique et lui sont
indifférentes.

\subsection{Théorème 4.2 : Théorème de découplage géométrique (forme Ruelle)}
\label{ssec:thm_4_2_fr}

\begin{theoreme}[Théorème de découplage géométrique, forme Ruelle]
\label{thm:ruelle_zero_fr}
\footnote{Vérifié formellement en Lean~4 + Mathlib comme
\texttt{PT.CrtDecoupling.Main.geometric\_decoupling\_ruelle} ;
cf.~\cite[\S~8.5]{PT_CRT_DECOUPLING}.}
Sous la mesure idéalisée $\pi_m^{\mathrm{Ruelle}}$, on a
$I(\Phi_p\,;\,\Phi_q) = 0$ exactement.
\end{theoreme}

La démonstration est immédiate à partir de la factorisation
tensorielle CRT du transfert
$T_{pq} = T_p \otimes T_q$ (Théorème~3.1
de~\cite{PT_CRT_DECOUPLING}) et de la factorisation correspondante
de l'unique vecteur propre gauche de Perron
$\pi_{pq}^{\rm Ruelle} = \pi_p^{\rm Ruelle} \otimes \pi_q^{\rm Ruelle}$
(Lemme~3.2 du même article). Cette factorisation est une
condition d'indépendance, donc la divergence de Kullback--Leibler
entre la loi jointe et le produit des marginales est nulle ;
par l'inégalité de traitement de données, l'information mutuelle
de toute paire d'observables bornées en hérite.

\emph{Validation numérique.} La factorisation
$\pi_{15}^{\rm Ruelle} = \pi_3^{\rm Ruelle} \otimes \pi_5^{\rm Ruelle}$
est vérifiée au sens du résidu
\begin{equation}
  \|\pi_{15}^{\rm direct} - \pi_3 \otimes \pi_5\|_2
  \;=\; 5{,}88 \times 10^{-53}
\end{equation}
à \texttt{mpmath} 50~chiffres de précision (test T2 PASS de
\texttt{scripts/test\_crt\_decoupling.py}~\cite{PT_CRT_DECOUPLING}).

\subsection{Corollaire 4.3 : Forme fermée du découplage}
\label{ssec:cor_4_3_fr}

\begin{corollaire}[Forme fermée du découplage]
\label{cor:closed_form_fr}
\footnote{Non formalisable comme théorème Lean en l'état car le
paramètre de profondeur $k_{p,q}^{\rm eff}$ est ajusté
empiriquement ; seule la prédiction à $k$ fixé admet un énoncé
Lean-vérifiable ;
cf.~\cite[\S~8.5]{PT_CRT_DECOUPLING}.}
La constante $\mathcal{I}_{p,q}$ admet une expression fermée
dérivée des angles d'holonomie :
\begin{equation}
  \mathcal{I}_{p,q}
  \;=\;
  \frac{1}{2 \ln 2}\Bigl[
    \cos^{2 k_{p,q}}(\theta_3, \mustar)
    + \sin^{4 k_{p,q}}(\theta_2, p_1)
  \Bigr].
  \label{eq:closed_form_fr}
\end{equation}
Validations numériques (cf.~\cite{PT_CRT_DECOUPLING}, Tableau B) :
$\mathcal{I}_{3,5} = 0{,}132$ bits (cible $0{,}138$, écart $4{,}5\%$) ;
$\mathcal{I}_{3,7} = 0{,}306$ bits (cible $0{,}283$, écart $8{,}1\%$).
\end{corollaire}

Le paramètre de profondeur $k_{p,q}$ est ajusté empiriquement à
$k_{3,5}^{\rm eff} = 7$ et $k_{3,7}^{\rm eff} = 4$ ; une dérivation
combinatoire à partir de la structure antidiagonale de $T_3$ reste
ouverte (cf.\ remarque~7.1 de~\cite{PT_CRT_DECOUPLING}). Cette
incertitude ne porte que sur la \emph{valeur quantitative} de
$\mathcal{I}_{p,q}$ ; le résultat structurel
(Théorème~\ref{thm:invariance_geom}) ne dépend pas du choix de
$k_{p,q}$.

\subsection{La non-localité spatiale dissoute en non-localité modulaire}
\label{ssec:dissolution}

Les trois énoncés précédents permettent de reformuler en termes
PT le statut de la « non-localité » apparente des corrélations
quantiques. On peut résumer la situation en trois propositions
strictement parallèles aux trois positions du débat fondationnel
(\S~\ref{sec:debat}), mais avec un statut différent :

\paragraph{(i) Les corrélations inter-canaux existent.}
$\mathcal{I}_{3,5}$ et $\mathcal{I}_{3,7}$ sont
non nuls et de magnitude prédite par le
Corollaire~\ref{cor:closed_form_fr}. La PT ne nie pas l'existence
des corrélations EPR/Bell ; elle en fournit, au contraire, une
prédiction quantitative.

\paragraph{(ii) Ces corrélations sont d'origine arithmétique.}
La forme fermée du Corollaire~\ref{cor:closed_form_fr}
fait intervenir uniquement les angles d'holonomie $\theta_3,
\theta_2$ évalués aux échelles canoniques $\mustar = 15$ et
$p_1 = 3$. Aucune intégrale spatiale, aucune métrique, aucune
distance entre points de l'espace n'entre dans l'expression. La
non-localité observée n'est pas une non-localité dans
l'espace : c'est une corrélation intra-CRT entre canaux modulaires
différents.

\paragraph{(iii) Ces corrélations sont invariantes par $G(\mu)$.}
Le Théorème~\ref{thm:invariance_geom} interdit que la géométrie
émergente \emph{couple} les canaux modulaires. Quelle que soit la
position dans la famille géométrique
$\mu \in [\mu_c, \infty)$ --- y compris à toute distance entre
détecteurs dans un dispositif Bell --- l'information mutuelle entre
canaux distincts reste la même.

Ces trois propositions ensemble dissolvent la non-localité
spatiale au sens fort de la position B : les corrélations
quantiques ne traversent pas l'espace parce qu'elles ne sont pas
dans l'espace. Elles sont dans la structure CRT, qui se manifeste
dans la géométrie émergente $G$ par projection, sans la traverser
ni l'utiliser comme support causal. Le paradoxe apparent d'une
information « plus rapide que $c$ » dans les tests de Bell est un
artéfact de la lecture spatiale d'une corrélation qui, dans la
substrat PT, n'est pas spatiale. Nous revenons sur les conséquences
expérimentales de cette dissolution dans
la~\S~\ref{ssec:bell_consequence}.
